% !TEX root = sga4.5-en.tex
% ENGLISH EDITION -- TRANSLATED FROM FRENCH (AI-assisted, needs human review).

\chapter{Report on the trace formula}\label{II}

In this text, I have tried to expound as directly as possible the
cohomological theory of Grothendieck of $L$-functions. I follow very closely
some of the lectures given by Grothendieck at the IHES in the spring of 1966. In
the spirit of this volume, no appeal will be made to SGA 5 -- save two
references to passages of Expos\'e XI, independent moreover of the rest of that
seminar.

In \S 10 (p.81-90) of \cite{de73}, the reader will find a summary of the
theory, in the crucial case of curves.










\section{Some preliminaries}\label{II:1}




\subsection{Notations}\label{II:1-1}

$p,q,\dF,\dF_q$: $p$ is a prime number, $q=p^f$ a power of $p$ and
$\dF$ an algebraic closure of the prime field $\dF_p$. For every power
$p^n$ of $p$, we denote by $\dF_{p^n}$ the subfield with $p^n$ elements of $\dF$.

$X_0,X$: We will often denote by a symbol with a subscript $0$ an object
over $\dF_q$. The same symbol, without $0$, then denotes the object deduced
from it by extension of scalars to $\dF$. For example, if $\sF_0$ is a sheaf
on a scheme $X_0$ over $\dF_q$, we denote by $\sF$ its inverse image on
$X=X_0\otimes_{\dF_q} \dF$.

$F$: If $X_0$ is a scheme over $\dF_q$, we call \emph{Frobenius
morphism} and denote by $F$ the endomorphism of $X_0$ which ``sends the point of
coordinates $x$ to that of coordinates $x^q$'': it is the identity on
the underlying topological space and, for $x$ a local section of the structure
sheaf, $F^* x = x^q$.

We still denote by $F$ the endomorphism of $X$ deduced from it by
extension of scalars. On the set $X(\dF) = X_0(\dF)$ of
rational points of $X$, $F$ acts as the Frobenius substitution
$\varphi\in\gal(\dF/\dF_q)$ defined by $\varphi(x)=x^q$. In particular,
$X^F=X_0(\dF_q)$.

In case of ambiguity, we will write $F_{(q)}$ rather than $F$.








\subsection{Frobenius correspondence}\label{II:1-2}

Let $X_0$ be a scheme over $\dF_q$, and $\sF_0$ a sheaf on $X_0$. I like
to see the \emph{Frobenius correspondence} (an $F$-endomorphism of
$\sF_0$) as follows:
\begin{equation*}\tag{1.2.1}\label{II:eq:1-2-1}
  F^* : F^*\sF_0 \iso \sF_0
\end{equation*}

We regard $\sF_0$ as the sheaf of local sections of a space
$[\sF_0]$ \'etale over $X_0$ ($[\sF_0]$ is an algebraic space in the sense of M.
Artin). Functoriality of $F$ provides a commutative diagram
\[\xymatrix{
  [\sF_0] \ar[r]^-F \ar[d]
    & [\sF_0] \ar[d] \\
  X_0 \ar[r]^F
    & X_0
}\]
Because $[\sF_0]$ is \'etale over $X_0$, this diagram is cartesian; it provides
an isomorphism $[\sF_0]\iso F^*[\sF_0]$ with inverse \eqref{II:eq:1-2-1}. We still
denote by $F^*$ the correspondences or morphisms deduced from $F$ by extension
of scalars or by functoriality, such as
$F^*:\h_c^\bullet(X,\sF)\to \h_c^\bullet(X,\sF)$.

For a formal definition, I refer to \cite[XI.1,2]{sga5}. Only
the case $q=p$ is considered there, but the general case is treated in the same way.




\subsection{}\label{II:1-3}

The Frobenius correspondence is functorial in $(X_0,\sF_0)$. If
$f_0:X_0\to Y_0$ is a morphism of $\dF_q$-schemes, $\sF_0$ a sheaf on
$X_0$, $\sG_0$ a sheaf on $Y_0$ and $u\in\hom(f_0^*\sG_0,\sF_0)$ an
$f_0$-morphism from $\sG_0$ to $\sF_0$, we have a commutative diagram of spaces
$f_0 F = F f_0$, and above it a commutative diagram of sheaves
\[\xymatrix{
  X_0 \ar[r]^-{f_0} \ar[d]^-F
    & Y_0 \ar[d]^-F \\
  X_0 \ar[r]^-{f_0}
    & Y_0
}\qquad\qquad
\xymatrix{
  \sF_0
    & \ar[l]_-u \sG_0 \\
  \sF_0 \ar[u]_-{F^*}
    & \ar[l] \sG_0 \ar[u]^-{F^*}
}\]
In particular, the correspondence on $f_{0*}\sF_0$ deduced from the Frobenius
correspondence of $(X_0,\sF_0)$ is the Frobenius correspondence of
$(U_0,f_{0*}\sF_0)$.




\subsection{}\label{II:1-4}

The natural morphism $Y\to Y_0$ is a limit of \'etale morphisms. It
follows that for every abelian sheaf $\sF_0$ on $X_0$, the inverse
images on $Y$ of the $\eR^if_0\sF_0$ are the $\eR^if_*\sF$.

On $\eR^if_*\sF$, we have
\begin{enumerate}[\indent a)]
  \item the correspondence
    $F^*\eR^i f_*\sF \iso \eR^i f_*\sF \xrightarrow{F^*} \eR^if_*\sF$ deduced
    by functoriality from the Frobenius correspondence;
  \item the correspondence $F^*\eR^if_*\sF\to\eR^if_*\sF$ inverse image on
    $Y$ of the Frobenius correspondence on $\eR^ii f_{0*}\sF_0$.
\end{enumerate}

We deduce from (\ref{II:1-3}) that these correspondences coincide. They
induce the initial term of a Frobenius action on the whole spectral
sequence
\[
  \h^p(Y,\eR^q f_*\sF) \rightarrow \h^{p+q}(X,\sF) \text{.}
\]
Similarly in cohomology with proper support, and for the spectral sequence
\[
  \h_c^p(Y,\eR^q f_* \sF) \rightarrow \h_c^{p+q}(X,\sF)
\]




\subsection{Change of finite field}\label{II:1-5}

Let $(X_0,\sF_0)$ over $\dF_q$, $n$ an integer, and $(X_1,\sF_1)$ over
$\dF_{q^n}$ deduced from $(X_0,\sF_0)$ by extension of scalars. We have
on the one hand the Frobenius correspondence $F_{(q)}:X_0\to X_0)$,
$F_{(q)}^*:F_{(q)}^*\sF_0\to \sF_0$, on the other hand
$F_{(q^n)}:X_1\to X_1$, $F_{(q^n)}^*:F_{(q^n)}^*\sF_1\to\sF_1$. One checks
easily that $F_{(q^n)}$ is deduced from the $n$-th power of $F_{(q)}$
by extension of scalars from $\dF_q$ to $\dF_{q^n}$; after extension of
scalars to $\dF$, we still have the same identity $F_{(q^n)}=F_{(q)}^n$
between endomorphisms of $X$ and correspondences on $X$.

If $x$ is a closed point of $X_0$, with $[k(x),\dF_q]=n$, and
$\bar x\in X_0(\dF) = X(\dF)$ is localized at $x$, $\bar x$ is a
fixed point of $F_{(q^n)}=F_{(q)}^n$. We denote by $F_x^*$ the endomorphism of
$\sF_{\bar x}$ induced by $F_{(q^n)}^*$. Up to isomorphism,
$(F_{\bar x}^*,\sF_{\bar x})$ does not depend on the choice of $\bar x$. The trace,
etc. of $F_x^*$ is denoted by a notation such as
$\tr(F_x^*,\sF)$.




\subsection{\texorpdfstring{$L$}{L}-function}\label{II:1-6}

Let $X_0$ be a scheme of finite type over $\dF_q$. We denote by $|X_0|$ the set of
its closed points and, for $x\in |X_0|$, we put $\deg(x)=[k(x),\dF_p]$. If
$\sF_0$ is a $\dQ_\ell$-sheaf on $X_0$ (for this notion, see
below) -- or a constructible flat sheaf of modules over a
commutative noetherian ring $A$, we denote by $L(X_0,\sF_0)$ the formal series
with constant term $1$, in $\dQ_\ell\llbracket t\rrbracket$ or in
$A\llbracket t\rrbracket$,
\[
  L(X_0,\sF_0) = \prod_{x\in|X_0|} \det\left(1-F_x^* t^{\deg(x)},\sF\right)^{-1} \text{.}
\]




\subsection{Remark}\label{II:1-7}

In definition (\ref{II:1-6}), we have $\deg(x)=[k(x),\dF_p]$ (absolute
degree) and not $\deg(x)=[k(x),\dF_q]$ (degree relative to $\dF_q$). This has
the effect that $L(X_0,\sF_0)$ depends only on the scheme $X_0$ and the sheaf
$\sF_0$, not on $q$ and the structural morphism $X_0\to\spec(\dF_q)$.
The indeterminate $t$ is an avatar of $p^{-s}$.




\subsection{The Galois point of view}\label{II:1-8}

This n$^\circ$ will not be used in the sequel of the report. For $\sF_0$ a
abelian sheaf on $X_0$, the Galois group $\gal(\dF/\dF_q)$ acts on
$\dF$, and by transport of structure on $\h_c^i(X,\sF)$. In particular, the
Frobenius substitution $\varphi$ (\ref{II:1-1}) induces an automorphism,
still denoted $\varphi$, of $\h_c^i(X,\sF)$. The Frobenius correspondence
for its part defines an endomorphism $F^*$ of $\h_c^i(X,\sF)$. We have
\begin{equation*}\tag{1.8.1}\label{II:eq:1-8-1}
  {F^*}^{-1} = \varphi \qquad \text{(on $\h_c^i(X,\sF)$).}
\end{equation*}

To see this, apply (\ref{II:1-4}) to the case where $Y_0$ is a point:
$Y_0=\spec(\dF_q)$. We find that $\h_c^i(X,\sF)$ is the fibre, at the
geometric point $\spec(\dF)$ of $Y_0$, of $\eR^i f_! \sF_0$, and that $F^*$ is
the fibre of the Frobenius correspondence of $\eR^i f_!\sF_0$. Let
$[\eR^i f_! \sF_0]$ as in (\ref{II:1-2}). We have
$\h_c^i(X,\sF) = [\eR^i f_!\sF_0](\dF)$ and
\begin{enumerate}[\indent a)]
  \item $\varphi$ acts by transport of structure, via its action on $\dF$,
  \item $F^*$ is the inverse of
    $F:[\eR^i f_!\sF_0](\dF) \to [\eR^i f_!\sF_0](\dF)$.
\end{enumerate}

We conclude by the identity $F=\varphi$ of (\ref{II:1-1}).

One of the interests of the Galois point of view is that it allows reasoning by
transport of structure.










\section{\texorpdfstring{$\dQ_\ell$}{Ql}-sheaves}\label{II:2}

The notion of $\dQ_\ell$-sheaf is developed in detail in
\cite[V,VI]{sga5}. We will only use the definitions and results
below.




\begin{definition_}\label{II:2-1}
Let $X$ be a noetherian scheme. A $\dZ_\ell$-sheaf $\sF$ on $X$ is a
projective system of sheaves $\sF_n$ ($n\geqslant 0$), with $\sF_n$ a
constructible sheaf of $\dZ/\ell^{n+1}$-modules \cite[IX.2]{sga4}, and such that
the transition morphism $\sF_n\to\sF_{n-1}$ factors through an isomorphism
\[
  \sF_n\otimes_{\dZ/\ell^{n+1}} \dZ/\ell^n \iso \sF_{n-1}\text{.}
\]
We say that $\sF$ is \emph{smooth} if the $\sF_n$ are locally constant.
\end{definition_}

The terminology of \cite{sga5} is different, there one says ``$\dZ_\ell$-constructible
sheaf'' for ``$\dZ_\ell$-sheaf,'' and ``twisted constant
constructible'' for ``smooth.''

Similarly, below, for $\dQ_\ell$-sheaves.




\subsection{}\label{II:2-2}

If $k$ is a separably closed field, a sheaf of $\dZ/\ell^n$-modules
on $\spec(k)$ identifies with a $\dZ/\ell^n$-module, and the functor
$\sF\mapsto \varprojlim \sF_n$ identifies the $\dZ_\ell$-sheaves on
$\spec(k)$ with $\dZ_\ell$-modules of finite type. This justifies the following
definition: the \emph{fibre} of the $\dZ_\ell$-sheaf $\sF$ at the
geometric point $x$ of $X$ is the $\dZ_\ell$-module $\varprojlim (\sF_n)_x$.




\subsection{}\label{II:2-3}

Suppose $X$ connected, and let $x$ be a geometric point of $X$. It is known that
the functor ``fibre at $x$'' is an equivalence of the category of
locally constant constructible sheaves of $\dZ/\ell^n$-modules with the
category of $\dZ/\ell^n$-modules of finite type equipped with a continuous
action of $\pi_1(X,x)$. By passage to the limit, we deduce the following
proposition.




\begin{proposition_}\label{II:2-4}
Under the hypotheses above, the functor ``fibre at $x$'' is an
equivalence of the category of smooth $\dZ_\ell$-sheaves on $X$ with
that of $\dZ_\ell$-modules of finite type equipped with a continuous action of
$\pi_1(X,x)$.
\end{proposition_}




\begin{proposition_}\label{II:2-5}
Let $\sF$ be a $\dZ_\ell$-sheaf on a noetherian scheme $X$. There exists
a finite partition of $X$ into locally closed subsets $X_i$ such that
the $\sF|X_i$ are smooth.
\end{proposition_}

Put $\gr_\ell^n\sF=\ell^n\sF_m/\ell^{m+1}\sF_m$ for $m\geqslant n$, and
$\gr_\ell\sF=\bigoplus \gr_\ell^n\sF$. It is a sheaf of modules over
the noetherian ring $\gr_\ell\dZ=\dZ/\ell[t]$, graded for the
$\ell$-adic filtration. It is a quotient of
$\gr_\ell^0\sF \otimes_{\dZ/\ell} \dZ/\ell[t]$, hence a constructible sheaf of
$\gr_\ell\dZ$-modules \cite[IX.2]{sga4}. and there exists a finite partition
of $X$ into locally closed subsets $X_i$, such that
$\gr_\ell \sF|X_i$ are locally constant. Each of the sheaves
$\gr_\ell^n\sF|X_i$ is then locally constant as well as the $\sF_m|X_i$,
which are successive extensions of them.




\subsection{}\label{II:2-6}

If $\sF$ and $\sG$ are two $\dZ_\ell$-sheaves on $X$, we put
$\hom(\sF,\sG) = \varprojlim(\sF_n,\sG_n)$; it is a $\dZ_\ell$-module.
Since $\hom(\sF_n,\sG_n) = \hom(\sF_m,\sG_n)$ for $m\geqslant n$, we still have
\[
  \hom(\sF,\sG) = \varprojlim_n \varinjlim_m \hom(\sF_m,\sG_n) \text{,}
\]
so that the functor $\sF\mapsto \text{``}\varprojlim\text{''}\sF_n$ is
fully faithful, from the category of $\dZ_\ell$-sheaves into that of
pro-sheaves on $X$ ($=$ pro-objects of the category of sheaves on $X$).
We will use it to identify the $\dZ_\ell$-sheaves with certain
pro-sheaves (those such that each $\sF_n=\sF/\ell^{n+1}\sF$ is a sheaf,
and that $\sF\iso \text{``}\varprojlim\text{''}\sF_n$). One checks easily
that if $\sF+\text{``}\varprojlim\text{''}\sF_n$, with $\sF_n$ a
constructible sheaf of $\dZ/\ell^n$-modules, for $\sF$ to be a
$\dZ_\ell$-sheaf, it is necessary and sufficient that the projective systems (in $n$)
$\sF_n/\ell^{k+1}\sF_n$ be essentially constant: the projective
system of the $\sF_k'=\varprojlim_n\sF_n/\ell^{k+1}\sF_n$ then verifies the
conditions of (\ref{II:2-1}), and $\sF=\text{``}\varprojlim\text{''}\sF_k'$.

In what follows, we abandon the notation $\sF_n$ of (\ref{II:2-1}); if
$\sF$ is a $\dZ_\ell$-sheaf, the sheaf $\sF_n$ of (\ref{II:2-1}) will be
denoted $\sF\otimes\dZ/\ell^{n+1}$.




\subsection{}\label{II:2-7}

In the abelian category of pro-sheaves, the subcategory of
$\dZ_\ell$-sheaves is stable under kernel, cokernel and extension. This is clear
for cokernels. For kernels, if $f:\sF\to\sG$ is a morphism of
$\dZ_\ell$-sheaves, one deduces easily from (\ref{II:2-4}, \ref{II:2-5})
and Artin-Rees that the projective system of the
$\ker(f:\sF\otimes\dZ/\ell^n\to\sG\otimes\dZ/\ell^n)$ verifies the criterion
given in (\ref{II:2-6}) above. There even exists an integer $r$ such that for
all $n$
\[
  \ker(f)\otimes\dZ/\ell^n = \ker(f:\sF\otimes\dZ/\ell^{n+r}\to\sG\otimes\dZ/\ell^{n+r})\otimes\dZ/\ell^n \text{.}
\]
The case of extensions is left to the reader.




\subsection{}\label{II:2-8}

A $\dZ_\ell$-sheaf $\sF$ is \emph{torsion free} if
$\ker(\ell:\sF\to\sF) = 0$. Each $\sF\otimes\dZ/\ell^n$ is then flat over
$\dZ/\ell^n$ (i.e. with fibres finite free $\dZ/\ell^n$-modules). One
deduces from (\ref{II:2-4}, \ref{II:2-5}) that for every $\dZ_\ell$-sheaf
$\sF$ there exists an integer $n$ such that $\sF/\ker(\ell^n)$ is torsion free.




\subsection{}\label{II:2-9}

The abelian category of \emph{$\dQ_\ell$-sheaves} on $X$ is the
quotient of that of $\dZ_\ell$-sheaves by the subcategory of
torsion $\dZ_\ell$-sheaves. In equivalent terms, its objects are the
$\dZ_\ell$-sheaves and, denoting by $\sF\otimes\dZ_\ell$ the $\dZ_\ell$-sheaf
$\sF$, seen as an object of the category of $\dQ_\ell$-sheaves, we have
\[
  \hom(\sF\otimes\dQ_\ell,\sG\otimes\dQ_\ell) = \hom(\sF,\sG)\otimes_{\dZ_\ell}\dQ_\ell \text{.}
\]

The \emph{fibre} of a $\dQ_\ell$-sheaf $\sF\otimes\dQ_\ell$ at a
geometric point $x$ of $X$ is the finite dimensional
$\dQ_\ell$-vector space $\sF_x\otimes_{\dZ_\ell}\dQ_\ell$.




\subsection{}\label{II:2-10}

Let $X$ be a separated scheme of finite type over an algebraically
closed field $k$ and $\sF$ a constructible $\dQ_\ell$-sheaf on $X$. Let $\sF'$ be a
constructible $\dZ_\ell$-sheaf such that
$\sF\sim \sF'\otimes_{\dZ_\ell}\dQ_\ell$. We will see in (\ref{II:4-11}) that
\[
  \left(\varprojlim \h_c^q(X,\sF'\otimes \dZ/\ell^n)\right)\otimes_{\dZ_\ell}\dQ_\ell
\]
is a finite dimensional $\dQ_\ell$-vector space. It depends only on
$\sF'$, and is denoted $\h_c^q(X,\sF)$. The proof does not use that $\ell$ is
prime to the characteristic, but only this case interests us.




\subsection{}\label{II:2-11}

In this expos\'e, all important computations will be done at the level of
$\dZ/\ell^n$-sheaves, passage to $\dQ_\ell$-sheaves being done only at the
very end. In practice, it is however often convenient to work
systematically with $\dQ_\ell$-sheaves. In Expos\'es
V and VI of \cite{sga5}, Jouanolou gives the formalism which makes this possible.
One of the essential results (VI 2.2) is that if $f:X\to Y$ is a
separated morphism of finite type of noetherian schemes, and $\sF$ a
$\dZ_\ell$-sheaf on $Y$, then, for each $q$, the pro-sheaf
$\text{``}\varprojlim\text{''}\eR^q f_! (\sF\otimes\dZ/\ell^n)$ is a
$\dZ_\ell$-sheaf on $Y$. He even proves a stronger
regularity property (in the style of Artin-Rees) for the projective system of the
$\eR^q f_!(\sF\otimes\dZ/\ell^n)$.

Jouanolou even places himself in a more general framework, encompassing the
very useful case of $E_\lambda$-sheaves, for $E_\lambda$ a finite extension of
$\dQ_\ell$. They are defined as the $\dQ_\ell$-sheaves were,
with $\dQ_\ell$ replaced by $E_\lambda$, $\dZ_\ell$ by the ring
$\sO_\lambda$ of the valuation $\lambda$ and $\dZ/\ell^n$ by
$\sO_\lambda/\pi^n$, for $\pi$ a uniformizer. It would amount to the same to
define an $E_\lambda$-sheaf as a $\dQ_\ell$-sheaf $\sF$, equipped
with a homomorphism $E_\lambda\to\operatorname{End}(\sF)$.










\section{Trace formula and \texorpdfstring{$L$}{L}-functions}\label{II:3}

\emph{We use the notations of \S\ref{II:1}, and $\ell$ is a prime
$\ne p$.}

Let $X_0$ be a separated scheme of finite type over $\dF_q$ ($q=p^f$) and
$\sF_0$ a constructible $\dQ_\ell$-sheaf on $X_0$. Grothendieck's
cohomological interpretation of $L$-functions is the following theorem:




\begin{theorem_}\label{II:3-1}
$L(X_0,\sF_0) = \prod_i \det\left(1-F^* t^f, \h_c^i(X,\sF)\right)^{(-1)^{i+1}}$.
\end{theorem_}

It is obtained as a consequence of the trace formula:




\begin{theorem_}\label{II:3-2}
For every integer $n$,
$\sum_{x\in X^{F^n}} \tr\left({F^n}^*,\sF_x\right) = \sum_i (-1)^i \tr\left({F^n}^*,\h_c^i(X,\sF)\right)$.
\end{theorem_}




Let us briefly recall the classical method to deduce (\ref{II:3-1})
from (\ref{II:3-2}). Put $T=t^f$; both sides are formal series
in $T$, with constant term $1$. Since $\dQ_\ell$ is of characteristic
$0$, it suffices to prove that the logarithmic derivatives
$T\frac{d}{dT}\log$ of the two sides are equal.

To compute these logarithmic derivatives, we will use the formula
\[
  T \frac{d}{dT}\log\det(1-f T)^{-1} = \sum_{n>0} \tr(f^n) T^n \text{,}
\]
% NOTE: the following lines are not typesetting correctly
whose proof (in a slightly more general framework) is recalled
below. For the first side, one finds
\begin{align*}
  T\frac{d}{dT}\log L(X_0,\sF_0)
    &= \sum_{x\in |X_0|} \sum_n [k(x):\dF_q] \tr\left(F_x^n,\sF\right) T^{n\cdot[k(x):\dF_q]} \\
    &= \sum_{n>0} T^n \sum_{x\in X^{F^n}} \tr\left({F^n}^*,\sF_x\right)
\end{align*}
and for the second side
\[
  \sum_{n>0} T^n \sum (-1)^i \tr\left({F^n}^*,\h_c^i(X,\sF)\right) \text{ ; }
\]
one compares term by term.




\begin{proposition_}\label{II:3-3}
Let $A$ be a commutative ring, $M$ a finite projective $A$-module and
$f$ an endomorphism of $M$. For every invertible formal series $Q$, put
$\frac{d}{dT}\log Q = Q^{-1}\frac{d}{dT}Q$. We have
\begin{equation*}\tag{3.3.1}\label{II:eq:3-3-1}
  T\frac{d}{dT}\log\det(1-f T)^{-1} = \sum_{n>0} \tr(f^n) T^n \text{.}
\end{equation*}
\end{proposition_}

Let $M'$ be an $A$-module such that $M\oplus M'$ is free of finite rank.
Replacing $(M,f)$ by $(M\oplus M',f\oplus 0)$ does not change the two sides
of (\ref{II:3-3}), and reduces us to the case where $M$ is free. For $M=A^d$,
\eqref{II:eq:3-3-1} is an algebraic identity concerning the coefficients
$f_i^j$ of $f$. By the principle of extension of algebraic identities,
it suffices to verify it for $A$ an algebraically closed field of
characteristic $0$. Both sides being additive in $M$ in an
extension $0\to M'\to M\to M''\to 0$, it even suffices to treat only the case
where $M$ is of rank $1$. If $f$ is multiplication by $a$,
\eqref{II:eq:3-3-1} then reduces to the formula
\[
  T\frac{d}{dT}\log\left((1-a T)^{-1}\right) = \sum_{n>0} a^n T^n \text{.}
\]




\begin{corollary_}\label{II:3-4}
Let $n$ be an integer, and $f^{(n)}$ the endomorphism
$(x_1,\dotsc,x_n)\mapsto \left(f(x_n),x_1,\dotsc,x_{n-1}\right)$ of $M^n$.
We have
\begin{equation*}\tag{3.4.1}\label{II:eq:3-4-1}
  \det(1-f T^n) = \det(1-f^{(n)} T) \text{.}
\end{equation*}
\end{corollary_}

Proceeding as in (\ref{II:3-3}), we reduce to assuming that $A$ is of
characteristic $0$. It then suffices to check that the logarithmic
derivatives $-T\frac{d}{dT}\log$ of the two sides are equal.
\begin{align*}
  -T\frac{d}{dT}\log\det(1-f T^n) &= - n T^n \frac{d}{dT^n} \log\det(1-f T^n) = n \sum_{m>0} \tr(f^m) T^{n m} \text{,} \\
  -T \frac{d}{dT}\log\det\left(1-f^{(n)} T\right) &= \sum_{m>0} \tr\left(f^{(n) m}\right) T^m \text{,}
\end{align*}
and one observes that
\begin{align*}
  \tr(f^{(n) m}) &= 0 \qquad \text{if $n\nmid m$} \\
  \tr(f^{(n) n m}) &= n \tr(f^m) \text{.}
\end{align*}




\subsection{Remark}\label{II:3-5}

For $A=\dQ_\ell$, \eqref{II:eq:3-4-1} is the particular case
$X_0=\spec(\dF_{q^n})$ of (\ref{II:3-1}). Using this identity, it is
easy to check directly that the second side of (\ref{II:3-1}) is
independent of $q$, and of the structural morphism $X_0\to\spec(\dF_q)$.




\subsection{Remark}\label{II:3-6}

If $j:X_0\hookrightarrow \bar X_0$ is a compactification of $X_0$, we have
$L(X_0,\sF_0) = L(\bar X_0, j_!\sF_0)$ and
$\h_c^i(X,\sF) = \h^i(\bar X,j_!\sF)$. This reduces (\ref{II:3-1}) to the
proper case, and explains the use of cohomology with proper supports. For $X$
proper, (\ref{II:3-2}) has the form of a Lefschetz trace formula. Note
that the local terms $\tr(F^{n *},\sF_x)$ ($x\in X^{F^n}$) do not
depend except on the fibre of $\sF$ at $x$, and are not affected with
multiplicity. For $X$ and $\sF$ smooth, this corresponds to the fact that the graph
of $F$ is transverse to the diagonal (since $d F = 0$).




\subsection{Remark}\label{II:3-7}

If one admits the formalism of $\dQ_\ell$-sheaves (cf. \ref{II:2-11}; one needs
the Leray spectral sequence and the base change theorem for the higher direct images
$\eR^q f_!$), it is easy to reduce the proof of (\ref{II:3-1}, \ref{II:3-2}), to the case where $X_0$ is
a smooth curve and where $\sF_0$ is smooth. This is clearly explained
in \cite[\S 5]{gr95} (for \ref{II:3-1}; \ref{II:3-2} is treated similarly).




\subsection{}\label{II:3-8}

We have two methods to prove (\ref{II:3-2}).


\paragraph{Lefschetz-Verdier}
If $X_0$ is proper (cf. \ref{II:3-6}), the general trace formula of
Lefschetz-Verdier allows one to express the second side of (\ref{II:3-2}) as
a sum of local terms, one for each point of $X^{F^n}$. In the original
version of \cite{sga5}, this formula was only proved modulo
resolution of singularities. The reader will find an unconditional proof
in the definitive version. In the case of curves, to which one can reduce
(\ref{II:3-7}), the ingredients of the proof were
moreover all available.

To deduce (\ref{II:3-2}) from the Lefschetz-Verdier formula, one must
be able to compute its local terms. For a curve and the
Frobenius endomorphism, this had been done by Artin and Verdier \cite{ve67}.


\paragraph{Nielsen-Wecken}
A method inspired by the work of Nielsen and Wecken allows one to reduce
(\ref{II:3-2}), to a particular case proved by Weil; this is what will be
explained in the following paragraphs.










\section{Reduction to theorems in \texorpdfstring{$\dZ/\ell^n$}{Z/l n}-cohomology}\label{II:4}

We will prove (\ref{II:3-2}) by passage to the limit from an
analogous theorem in $\dZ/\ell^n$-cohomology. Stating it presents the
following difficulty. Let $X_0$ be a separated scheme of finite type over
$\dF_q$, and $\sF_0$ a flat constructible sheaf of $\dZ/\ell^n$-modules.
Its fibres are finite free $\dZ/\ell^n$-modules, and the left hand
side of (\ref{II:3-2}) therefore makes sense (it is an element of $\dZ/\ell^n$).
On the other hand, the $\h_c^i(X,\sF)$ will in general not be free, so that
the right hand side is not defined. Overcoming this difficulty is
one of the essential uses made by Grothendieck of the theory of
derived categories.




\subsection{}\label{II:4-1}

For the definition of the derived category $\eD(\cA)$ of an
abelian category $\cA$, and that of the subcategories $\eD^+(\cA)$, $\eD^-(\cA)$,
$\eD^b(\cA)$, I refer to Verdier's text in this volume. For $\Lambda$
a ring, we write $\eD(\Lambda)$ instead of
$\eD(\text{category of left $\Lambda$-modules})$; for $X$ a site, we
write $\eD(X,\Lambda)$ instead of
$\eD($category of sheaves of left $\Lambda$-modules$)$. Similarly
for $\eD^-$, etc\ldots

The sign $\eR$ (resp. $\mathsf{L}$) indicates a right (resp.
left) derived functor. For example, $\lotimes$ indicates the derived functor of the tensor
product: for $K$ in $\eD^-(\Lambda^\circ)$ and $L$ in $\eD^-(\Lambda)$
($\Lambda^\circ$ ring opposite to $\Lambda$), $K\lotimes L$ is computed
as follows: one takes quasi-isomorphisms $K'\to K$, $L'\to L$, with
$K'$ and $L$ bounded above and one of them with flat components, and
the simple complex $s(K'\otimes_\Lambda L)$ associated to the double complex
$K'\otimes_\Lambda L'$.




\subsection{}\label{II:4-2}

Even to treat only the case of $\dZ/\ell^n$-cohomology, we will
need the theory of traces of endomorphisms of modules over
not necessarily commutative rings (group algebras of finite groups
$\dZ/\ell^n[G]$). These traces were introduced by Stallings and Hattori
\cite{ba76}.

Let $\Lambda$ be a ring (with unit, but not necessarily commutative).
We denote by $\Lambda^\natural$ the quotient of the additive group of $\Lambda$ by the
subgroup generated by the commutators $a b - b a$. If $f$ is an
endomorphism of the free left $\Lambda$-module $\Lambda^n$, with matrix
$f_i^j$, we denote by $\tr(f)$ the image of $\sum f_i^f$ in $\Lambda^\natural$.
If $f$ and $g$ are two homomorphisms
$f:\Lambda^n \leftrightarrows \Lambda^m : g$, one checks at once that
$\tr(f g) = \tr(g f)$.

Let $f$ be an endomorphism of a finite projective $\Lambda$-module $P$.
Write $P$ as a direct factor of a free module $\Lambda^n$: this means
choosing a module $P'$, and an isomorphism $\alpha:P\oplus P'\iso \Lambda^n$,
i.e. choosing a homomorphism $a:P\to \Lambda^n$, and a retraction
$b:\Lambda^n\to P$ (take $P=\ker(b)$). Let
$f'=\alpha(f\oplus 0)\alpha^{-1} =af b$. We put $\tr(f) = \tr(f')$. This
element of $\Lambda^\natural$ depends only on $f$: if another choice
$c:P\leftrightarrows \Lambda^m:d$ is made, we have $a=a(dc)(ba)$, whence
\[
 \tr(a f b) = \tr(adcbafb) = \tr(cbafbad) = \tr(cfd) \text{.}
\]

If $f$ is an endomorphism of a $\dZ/2$-graded finite projective $\Lambda$-module,
with components $f_i^j:P^j\to P^i$, we put
\[
  \tr(f) = \tr(f_0^0) - \tr(f_1^1) \text{.}
\]
If there is a risk of ambiguity, we will write $\tr(f;P)$ rather. If the
homomorphisms $f:P\leftrightarrows Q:g$ are homogeneous, one checks
at once by reduction to the free case that
\begin{equation*}\tag{4.2.1}\label{II:eq:4-2-1}
  \tr(f g) = \tr\left((-1)^{\deg f \deg g} g f\right) \text{.}
\end{equation*}

If $f$ is an endomorphism of a filtered $\dZ/2$-graded $\Lambda$-module
$(P,F)$, with finite filtration, compatible with the grading, and the
$\gr_F^i(P)$ are finite projective, then $P$ is finite projective,
and
\begin{equation*}\tag{4.2.2}\label{II:eq:4-2-2}
  \tr(f;P) = \sum \tr\left(f; \gr_F^i(P) \right)
\end{equation*}
(split the filtration to reduce to the case of a direct sum).




\subsection{}\label{II:4-3}

A $\dZ$-graded module will be considered as $\dZ/2$-graded by the
parity of the degree. In particular, if $f$ is an endomorphism of a bounded
complex of finite projective $\Lambda$-modules, we put
\begin{equation*}\tag{4.3.1}\label{II:eq:4-3-1}
  \tr(f) = \sum (-1)^i \tr(f^i) \text{.}
\end{equation*}

By \eqref{II:eq:4-2-1}, if $f$ is homotopic to zero, $\tr(f) = 0$:
\begin{equation*}\tag{4.3.2}\label{II:eq:4-3-2}
  f = d H + h D \rightarrow \tr(f) = 0 \text{.}
\end{equation*}

Let $\mathsf{K}_\text{parf}(\Lambda)$ be the category with objects bounded
complexes of finite projective $\Lambda$-modules, and arrows
morphisms of complexes taken up to homotopy. The functor
$\mathsf{K}_\text{parf}(\Lambda)\to\eD(\Lambda)$ is fully faithful. We denote
by $\mathsf{D}_\text{parf}(\Lambda)$ its essential image. Formulas
\eqref{II:eq:4-3-1} \eqref{II:eq:4-3-2} allow one to define $\tr(f)$ for $f$
an endomorphism of $K\in\ob \mathsf{D}_\text{parf}(\Lambda)$ (and this trace
depends only on the isomorphism class of $(K,f)$).




\subsection{}\label{II:4-4}

Recall that the filtered derived category $\mathsf{DF}(\cA)$ of an abelian
category $\cA$ is the category deduced from that of complexes of
filtered objects with finite filtration of $\cA$ by inverting filtered
quasi-isomorphisms ($=$ morphisms of filtered complexes inducing a quasi-isomorphism
on the associated graded). We have functors ``underlying complex'':
$\mathsf{DF}(\cA)\to\eD(\cA): (K,F)\mapsto K$ and ``$p$-th component of the
graded'': $\mathsf{DF}(\cA)\to\eD(\cA): (K,F)\mapsto \gr_F^p(K)$.

Let $\mathsf{KF}_\text{parf}(\Lambda)$ be the category with objects bounded
filtered complexes with finite filtration $(K,F)$ with $\gr_F^p(K^q)$ projective
of finite type, and arrows morphisms of complexes respecting the
filtration, taken up to homotopy respecting the filtration. The functor
$\mathsf{KF}_\text{parf}(\Lambda)\to\mathsf{DF}(\Lambda)$ is fully faithful;
its essential image $\mathsf{DF}_\text{parf}(\Lambda)$ is almost all zero;
if $(K,F)$ is in $\mathsf{D}_\text{parf}(\Lambda)$, then $K$ is in
$\mathsf{D}_\text{parf}(\Lambda)$.

By \eqref{II:eq:4-2-2}, if $f$ is an endomorphism of
$(K,F)\in\ob\mathsf{DF}_\text{parf}(\Lambda)$, we have
\begin{equation*}\tag{4.4.1}\label{II:eq:4-4-1}
 \tr(f,K) = \sum_p \tr(f,\gr_F^p(K)) \text{.}
\end{equation*}




\subsection{}\label{II:4-5}

Let $\Lambda$ be a ring, and $K\in\ob\eD^-(\Lambda)$. One says that $K$ is of
tor-dimension $\leqslant r$ if for every right $\Lambda$-module $N$, the
hypertors $\h^i(N\lotimes_\Lambda K)$ vanish for $i<-r$. It is said to be of finite
tor-dimension if this condition is verified for some $r$.

The same terminology is used for a complex of sheaves
$K\in\ob\eD^-(X,\Lambda)$.




\begin{lemma}\label{II:4-5-1}
Let $\Lambda$ be a left noetherian ring, and $K\in\ob\eD^-(\Lambda)$.
For $K$ to be in $\mathsf{DF}_\text{parf}(\Lambda)$, it is necessary and sufficient
that it be of finite tor-dimension and that the $\h^i(K)$ be of finite type.
\end{lemma}

Up to replacing $K$ by an isomorphic complex (in $\eD^-(\Lambda)$), one
can assume $K$ bounded above and with finite free components
(cf. \ref{II:4-7}, below). If $K$ is of tor-dimension $\leqslant r$,
the $\h^i(K)$ vanish for $i<-r$, $K^{-r}/\im(d)$ admits the free resolution
\[
  (\cdots \to K^{-r-1} \to K^{-r}) \to K^{-r}/\im(d) \text{,}
\]
and for every right $\Lambda$-module $N$,
$\h^{-r-k}(N\lotimes_\Lambda K) = \h^{-r-k}(N\otimes_\Lambda K)
= \tor_k(N,K^r/\im d)$ for $k\geqslant 1$.

By hypothesis, these groups vanish; the module $K^r/\im d$ is hence flat of
finite presentation, i.e. finite projective. The quotient
\[
  0 \to K^r/\im(d) \to K^{r+1} \to K^{r+2}\to \cdots
\]
of $K$ is quasi-isomorphic to $K$, and this shows that
$K\in\ob\mathsf{DF}_\text{parf}(\Lambda)$.




\begin{prop-def_}\label{II:4-6}
Let $X$ be a noetherian scheme, and $\Lambda$ a left noetherian ring.
We denote by $\eD_\textnormal{ctf}^b(X,\Lambda)$ the subcategory of
$\eD^-(X,\Lambda)$ formed of complexes $\sK$ satisfying the following
equivalent conditions:
\begin{enumerate}[\indent i)]
  \item $\sK$ is isomorphic (in $\eD^-(X,\Lambda)$) to a bounded complex of
    flat constructible sheaves of $\Lambda$-modules.
  \item $\sK$ is of finite tor-dimension and the sheaves $\h^i(K)$ are
    constructible.
\end{enumerate}
\end{prop-def_}

Same argument as in (\ref{II:4-5}), starting from the following lemma.




\begin{lemma_}\label{II:4-7}
If a complex $\sK$ of sheaves of $\Lambda$-modules is such that the
$\h^i(\sK))$ are constructible, and zero for large $i$, there exists a
quasi-isomorphism $\sK'\to \sK$, with $\sK'$ bounded above with
constructible flat components.
\end{lemma_}

We construct $\sK'$ proceeding from right to left. For $m$ such that the
$\h^i(\sK)$ ($i\geqslant m$) vanish, we take ${\sK'}^i = 0$ for
$i\geqslant m$. Suppose next already constructed the ${\sK'}^i$ for
$i>n$:
\[\xymatrix{
  \cdots \ar[r]
    & \sK^{n-1} \ar[r]
    & \sK^n \ar[r]
    & \sK^{n+1} \ar[r]
    & \cdots \\
  & & & {\sK'}^{n+1} \ar[u] \ar[r]
    & \cdots \text{,}
}\]
with $\h^i(\sK')\iso \h^i(\sK)$ for $i>n+1$, and $\ker(d)\to\h^{n+1}(\sK)$ surjective
in degree $n+1$. Define sheaves $\sA$ and $\sB$ by the cartesian diagram
\[\xymatrix{
  \sK^n \ar[r]
    & \sK^n/\im(d) \ar[r]
    & \ker(d) \\
  \sA \ar[u] \ar[r]^u
    & \sB \ar[u]\ar[r]
    & \ker(d) \ar[u] \text{.}
}\]

Then, $\sB$ is constructible, $u$ is surjective, and to advance one step, it
suffices to construct $K^n$ flat and constructible and $v:\sK^n\to \sA$ such that u
$u v$ is surjective. That this is possible is guaranteed by the following lemma.




\begin{lemma_}\label{II:4-8}
Let $v:\sA\to \sB$ be an epimorphism of sheaves of $\Lambda$-modules, with
$\sB$ constructible. There then exists $u:\sK\to \sA$ with $\sK$ flat and
constructible and $v u$ an epimorphism.
\end{lemma_}

The sheaf $\sA$ is a quotient of a sum of sheaves of the form
$\varphi_! \Lambda$, for $\varphi:U\to X$ \'etale. Since $\sB$ is noetherian
(\cite[IX.2.10]{sga4}), a finite sum of these already maps onto $\sB$,
whence the lemma.




\begin{theorem_}\label{II:4-9}
Let $f:X\to Y$ be a separated morphism of finite type of noetherian schemes.
If $\sK\in\ob\eD_\textnormal{ctf}^b(X,\Lambda)$, then
$\eR f_!\sK\in \ob\eD_\textnormal{ctf}^b(Y,\Lambda)$.
\end{theorem_}

Recall that if $j:X\hookrightarrow \bar X \xrightarrow{\bar f} Y$ is a
relative compactification of $X/Y$, we have by definition
$\eR f_! = \eR \bar f_* \circ j_!$, where $\eR \bar f_*$ is the derived functor
of the functor $\bar f_*$. This formula reduces the proof of (\ref{II:4-9})
to the case where $f$ is proper. For $f$ proper, $f_*$ is of finite
cohomological dimension; this allows one to define $\eR f_*$ on the whole
derived category, and also as a functor
\[
  \eR f_* : \eD^-(X,\Lambda) \to \eD^-(Y,\Lambda) \text{.}
\]
By composition, one likewise defines in general
\[
  \eR f_! : \eD^-(X,\Lambda) \to \eD^-(Y,\Lambda) \text{.}
\]

For every right $\Lambda$-module $N$, we have (proof given below)
\begin{equation*}\tag{4.9.1}\label{II:eq:4-9-1}
  N\lotimes \eR f_! \sK \iso \eR f_! (N\lotimes \sK) \text{.}
\end{equation*}

Deduce (\ref{II:4-9}) from \eqref{II:eq:4-9-1}.
\begin{enumerate}[\indent a)]
  \item The spectral sequence
    \[
      E_2^{p q} = \eR^p f_! \h^q(\sK) \rightarrow \h^{p+q} \eR f_! \sK \text{.}
    \]
    and the finiteness theorem for $\eR f_!$, show that
    $\eR f_!\sK\in \ob \eD^b(X,\Lambda)$ and is with constructible cohomology.
    It remains to prove finite tor-dimension (\ref{II:4-6}).
  \item If $\sK$ is of tor-dimension $\leqslant -r$, so is
    $\eR f_! \sK$: we have $\h^i(N\lotimes \eR f_! \sK) = \h^i(\eR f_!(N\lotimes \sK))$,
    and, in the above spectral sequence for $N\lotimes \sK$, the
    $E_2^{p q}$ vanish for $q<r-s$ a fortiori for $p+q<r$.
\end{enumerate}

Prove \eqref{II:eq:4-9-1}, staying at the level of complexes.
\begin{enumerate}[\indent a)]
  \item Using the defining formula $\eR f_! = \eR f_*\circ j_!$, we reduce
    to assuming $f$ proper.
  \item Represent $\sK$ by a complex bounded
    above with $f_*$-acyclic components: $\eR^p f_* \sK^q = 0$
    for $p>0$. We then have $\eR f_* K \sim f_* \sK$. It is possible to
    work thus with complexes bounded above only because
    $f_*$ is of finite cohomological dimension.
  \item If $L$ is a free right $\Lambda$-module, we have
    $\eR^p f_*(L\otimes \sK^q) = L\otimes \eR^p f_* \sK^q$: this is trivial for
    $\Lambda$ free of finite type, and the general case follows from it by
    passage to the limit. The sheaf $L\otimes \sK^q$ is acyclic for
    $f_*$, and $f_*(L\otimes \sK^q) = L\otimes f_* \sK^q$.
  \item Let $N_\bullet$ be a free resolution of $N$. We have
    $N\lotimes \sK\sim s(N_\bullet\otimes \sK)$ (simple complex associated to the
    double complex $N_\bullet\otimes \sK$). By c),
    \[
      \eR f_*(N\lotimes \sK) \sim \eR f_* (s(N_\bullet\otimes \sK)) \sim f_*(s(N_\bullet\otimes \sK)) \sim s(N_\bullet\otimes f_* \sK) \sim N\lotimes \eR f_* \sK \text{,}
    \]
    i.e. the announced formula.
\end{enumerate}

When $Y$ is the spectrum of a separably closed field, the functor $\eR f_!$
identifies with a functor denoted
$\eR \Gamma_c : \eD_\text{ctf}^b(X,\Lambda)\to\eD_\text{parf}(\Lambda)$.




\begin{theorem_}\label{II:4-10}
With the notations of \S\ref{II:1}, let $X_0$ be a separated scheme of
finite type over $\dF_q$, and $\Lambda$ a noetherian torsion ring, killed by an
integer prime to $p$. Let $\sK_0\in\ob\eD_\textnormal{ctf}^b(X_0,\Lambda)$.
With the definition (\ref{II:4-3}) of traces, we have for every $N$
\begin{equation*}\label{II:eq:4-10}
  \sum_{x\in X^{F^n}} \tr\left({F^n}^*,\sK_x\right) = \tr\left({F^n}^*,\eR\Gamma_c(X,\sK)\right) \text{.}
\end{equation*}
\end{theorem_}

In the end of this \S, we will deduce \ref{II:3-2} from \ref{II:4-10}, and
prove \ref{II:2-10}.




\subsection{Proof of \texorpdfstring{\ref{II:2-10}}{2.10}}\label{II:4-11}

Let $X$ be a separated scheme of finite type over an algebraically closed
field $k$, and let $\sG$ be a $\dQ_\ell$-sheaf on $X$. By \ref{II:2-8} one
can write it in the form $\sG=\sF\otimes\dQ_\ell$, with $\sF$ a
torsion free $\dZ_\ell$-sheaf. Put
\[
  K_n = \eR \Gamma_c(X,\sF\otimes\dZ/\ell^{n+1})\in\eD(\dZ/\ell^{n+1}) \text{.}
\]
By \ref{II:4-9}, we have $K_n\in\ob\eD_\text{parf}(\dZ/\ell^{n+1})$, and it
follows from \eqref{II:eq:4-9-1} that the natural arrow
\begin{equation*}\tag{4.11.1}\label{II:eq:4-11-1}
\xymatrix{
  K_n\lotimes_{\dZ/\ell^{n+1}} \dZ/\ell^n \ar[r] & K_{n+1} \text{,}
}
\end{equation*}
is an isomorphism. The problem of the definition of this arrow is
discussed in \ref{II:4-12}.

One can now invoke \cite[XI.3.3]{sga5} (p. 31, 1-2 at the end; this
passage is independent of the rest of \cite{sga5}) to realize each
$K_n$ by a bounded complex of finite free $\dZ/\ell^{n+1}$-modules,
and the arrows \eqref{II:eq:4-11-1} by \emph{isomorphisms of complexes}
\[\xymatrix{
  K_n\otimes \dZ/\ell^n \ar[r]^-\sim
    & K_{n-1} \text{.}
}\]
Let $K$ be the complex $\varprojlim K_n$. It is a bounded complex of
free $\dZ_\ell$-modules, and $K_n\simeq K\otimes_{\dZ_\ell}\dZ/\ell^{n+1}$. We
have
\[
  \h^i(K) = \varprojlim \h^i(K_n) \text{,}
\]
(every projective system of finite groups indeed verifies the
Mittag-Leffler condition), and $\h^i(K)$ is a $\dZ_\ell$-module of finite type. We have
\[
  \h_c^i(X,\sG) = \left(\varprojlim \h^i(K_n) \right)\otimes_{\dZ_\ell} \dQ_\ell \text{:}
\]
it is a finite dimensional $\dQ_\ell$-vector space.




\subsection{The arrow \texorpdfstring{\eqref{II:eq:4-11-1}}{(4.11.1)}}\label{II:4-12}

To \emph{define} the arrow \eqref{II:eq:4-11-1}, one cannot
proceed as in \ref{II:4-9}, i.e. resolve $\dZ/\ell^n$ by a complex
of free $\dZ/\ell^{n+1}$-modules, since one wants to construct an isomorphism
in $\eD_\text{parf}(\dZ/\ell^n)$. The desired arrow is a particular case of
the following base-change arrow.

(*) Let $\Lambda\to\Lambda'$ be a homomorphism of noetherian torsion rings.
For $X$ a separated scheme of finite type over $k$ algebraically
closed, and $\sK\in\ob\eD_\text{ctf}(X,\Lambda)$, we have an isomorphism
\[\xymatrix{
  \eR\Gamma_c(X,\sK)\lotimes_\Lambda \Lambda' \ar[r]^-\sim & \eR \Gamma_c(X,\sK\lotimes_\Lambda \Lambda') \text{.}
}\]

More general arrows are constructed in \cite[XVII 4.2.12]{sga4}. The
method is the following:
\begin{enumerate}[\indent a)]
  \item The definition of $\eR\Gamma_c$ and the formula
    $j_!(\sK\otimes_\Lambda\Lambda') = (j_! \sK)\otimes_\Lambda\Lambda'$ for an
    open immersion reduce us to the case where $X$ is proper.
  \item For $X$ proper, one replaces $\sK$ by a bounded complex with
    components acyclic for the functor $\Gamma$, and with fibre at every
    geometric point (or only at every closed point) homotopic to a
    complex of flat $\Lambda$-modules. For such a complex, we have
    $\Gamma(X,\sK)\sim \eR\Gamma(X,\sK)$ and
    $K\otimes_\Lambda\Lambda'\sim \sK\lotimes_\Lambda\Lambda'$; the arrow
    \eqref{II:eq:4-11-1} is defined as the composition
    \[
      \eR\Gamma(X,\sK)\lotimes_\Lambda \Lambda' \to \Gamma(X,\sK)\otimes_\Lambda\Lambda' \to \Gamma(X,\sK\otimes_\Lambda\Lambda') \xleftarrow\sim \eR\Gamma(X,\sK\lotimes_\Lambda\Lambda') \text{.}
    \]
  \item It remains to show that it is possible to replace $\sK$ by a
    complex of the desired type \cite[XVII.4.2.10]{sga4}. One first replaces $\sK$ by
    a bounded complex with flat components (\ref{II:4-6}), then one takes
    a truncated canonical flasque resolution of it (this does not modify the homotopy
    type of the fibre complexes, and one truncates far enough for
    the acyclicity condition for $\Gamma$ to be verified (i.e. beyond
    the cohomological dimension of $\Gamma$)).
\end{enumerate}




\subsection{Proof of \texorpdfstring{(\ref{II:3-2})}{(3.2)}}\label{II:4-13}

To simplify the notations, we will treat only the case $n=1$ of
(\ref{II:3-2}). The general case follows from it by replacing below
$F$ by $F^n$.

Let therefore $X_0$ be separated of finite type over $\dF_q$, and $\sG_0$ a
$\dQ_\ell$-sheaf on $X_0$. We have $\sG_0=\sF_0\otimes\dQ_\ell$, for
a suitable torsion free $\dZ_\ell$-sheaf $\sF_0$. Let
$K_n=\eR\Gamma_c(X,\sF\otimes\dZ/\ell^n)$. The Frobenius morphisms induce
morphisms
\[
  F^* : K_n\to K_n \qquad \text{(in $\eD_\text{parf}(\dZ/\ell^{n+1})$),}
\]
which are deduced from one another via the isomorphisms
\eqref{II:eq:4-11-1}.

Realize as in (\ref{II:4-11}) the $K_n$ and the isomorphisms
\eqref{II:eq:4-11-1} at the level of complexes. The already cited passage of
\cite[XV.3.3]{sga5} allows one to realize the morphisms $F^*$ by
endomorphisms of complexes, still denoted $F^*$, which reduce to one
another. We still denote by $F^*$ their projective limit, and the
endomorphisms deduced from it. Since
$\h_c^i(X,\sG) = \h^i(K)\otimes\dQ_\ell = \h^i(K\otimes\dQ_\ell)$, we have (with
the sign convention of (?.1)
\[
  \tr\left(F^*, \h_c^\bullet(X,\sG)\right) = \tr\left(F^*,K^\bullet\otimes\dQ_\ell\right) = \varprojlim_n \tr\left(F^*,K_n^\bullet\right) = \varprojlim_n \tr\left(F^*,\eR\Gamma_c(X,\sF\otimes\dZ/\ell^{n+1})\right) \text{.}
\]
Apply (\ref{II:4-10}) to $\sF_0\otimes\dZ/\ell^{n+1}$ (seen as a complex
concentrated in degree $0$), for $\Lambda=\dZ/\ell^{n+1}$. We find that
\begin{align*}
  \tr\left(F^*,\eR\Gamma_c(X,\sF\otimes\dZ/\ell^{n+1})\right)
    &= \sum_{x\in X^F} \tr\left(F^*,\sF_x\otimes\dZ/\ell^{n+1}\right) \\
    &= \sum_{x\in X^F} \tr\left(F^*,\sG_x\right) \mod \ell^{n+1} \text{,}
\end{align*}
and (\ref{II:3-2}), follows from it by passage to the limit.










\section{The Nielsen-Wecken method}\label{II:5}

In this chapter, we prove two particular cases of (\ref{II:4-10}). The
general case will be considered in the next chapter.




\begin{lemma_}\label{II:5-1}
Theorem (\ref{II:4-10}) is true for $X_0$ of dimension $0$.
\end{lemma_}

In dimension $0$, formula \eqref{II:eq:4-10} holds for every
correspondence, and not only for iterates of a Frobenius. It is
a particular case of the following statement (applied to $X(\dF)$, $K$ and
the ${F^*}^n$).

(*) Let $X$ be a finite set, $F:X\to X$, $\sK\in\ob\eD_\text{parf}(X,\Lambda)$
and $F^*:F^* \sK\to \sK$. We have
\[
  \sum_{X^F} \tr\left(F^*,\sK_x\right) = \tr\left(F^*,\Gamma(X,K)\right) \text{.}
\]

We have $\Gamma(X,\sK)=\bigoplus_{x\in X} \sK_x$, and the formula follows from the fact that
for every morphism $u:\bigoplus_{x\in X} \sK_x\to \bigoplus_{x\in X} \sK_x$, with
matrix $u_x^y$, we have $\tr(u)=\sum_{x\in X} \tr\left(u_x^x\right)$.




\begin{lemma_}\label{II:5-2}
Let $\Lambda$ be a noetherian torsion ring, killed by a power of the
prime $\ell\ne p$, $X_0$ a projective, smooth and connected curve over
$\dF_q$, $j:U_0\hookrightarrow X_0$ a dense open of $X_0$ and $\sG_0$ a
locally constant sheaf of projective $\Lambda$-modules on $U_0$. We have
\[
  \sum_{x\in U^F} \tr\left(F^*,\sG_x\right) = \tr\left(F^*,\eR\Gamma_c(U,\sG)\right) \text{.}
\]
\end{lemma_}

More precisely, we will deduce (\ref{II:5-2}) from the




\begin{theorem_}\label{II:5-3}
Let $X$ be a non-singular projective curve over an algebraically closed
field $k$, $f$ an endomorphism with isolated fixed points of $X$ and $v(f)$ the
number of fixed points of $f$, each counted with its multiplicity:
$v(f)$ is the intersection number of the graph of $f$ with the diagonal of
$X\times X$. Then, for $\ell$ prime to the characteristic exponent of $k$,
we have
\[
  v(f) = \sum (-1)^i \tr\left(f^*,\h^i(X,\dQ_\ell)\right) \text{.}
\]
\end{theorem_}

This theorem is due to Weil (\cite{we48}, n$^\circ$ 43, 66 and 68) and the
proof of Weil is reproduced in Lang \cite{la83} (VI \S3 combined
with VII \S2 th.3). It can also be deduced from the formalism of the
cohomology class associated to a cycle (\nameref{IV}, \ref{IV:3-7}).




\begin{corollary_}\label{II:5-4}
Let $j:U\hookrightarrow X$ be a dense open of $X$, with $f^{-1}(U)=U$. If
the fixed points of $f$ in $X\setminus U$ are of multiplicity one, then the
number $v_U(f)$ of fixed points of $f$ in $U$ (each counted with its
multiplicity) is
\[
  v_U(f) = \sum (-1)^i \tr\left(f^*,\h_c^i(U,\dQ_\ell)\right) \text{.}
\]
\end{corollary_}

Let $F$ be the finite set $X\setminus U$. The long exact cohomology sequence
defined by the short exact sequence
$0\to j_! \dQ_\ell\to\dQ_\ell\to{\dQ_\ell}_F\to 0$:
\[\xymatrix{
  \ar[r]^-\partial
    & \h_c^i(U,\dQ_\ell) \ar[r]
    & \h^i(X,\dQ_\ell) \ar[r]
    & \h^i(F,\dQ_\ell) \ar[r]^-\partial
    &
}\]
provides, with the sign convention of (\ref{II:3-1}).
\[
  \tr\left(f^*,\h_c^\bullet(U,\dQ_\ell)\right)
    = \tr\left(f^*,\h^\bullet(X,\dQ_\ell)\right) - \tr\left(f^*,\h^\bullet(F,\dQ_\ell)\right)
    = v(f) - \tr(f^*,\dQ_\ell^F) \text{.}
\]
The trace of $f^*$ on $\dQ_\ell^F$ is the number $v_F(f)$ of fixed points of
$f$ in $F$; the multiplicity-one hypothesis therefore ensures that
\[
  \tr\left(f^*,\h_c^\bullet(U,\dQ_\ell)\right)
    = v_F(f) - v_U(f) \text{,}
\]
as claimed.

The proof of (\ref{II:5-2}) uses a few lemmas which we will now
develop on a particular case of noncommutative traces
(\ref{II:4-2}). \emph{For brevity, we will say projective for finite projective}.




\subsection{}\label{II:5-5}

Let $\Lambda$ be a ring, and $H$ a finite group. The map
$\sum a_h \cdot a\mapsto $ class of $a_e$, from the group algebra
$\Lambda[H]$ to $\Lambda^\natural$, vanishes on the commutators $ab-ba$ of
$\Lambda[H]$. It hence factors through
\[
  \varepsilon : \Lambda[H]^\natural \to \Lambda^\natural \text{.}
\]
If $F$ is an endomorphism of a projective $\Lambda[H]$-module $P$, we put
\[
  \tr_\Lambda^H(F) = \varepsilon \tr_{\Lambda[H]} (F)
\]
where on the second side appears the trace (\ref{II:4-2}). To avoid
ambiguities, we will sometimes write this trace $\tr_\Lambda^H(F,P)$.




\begin{proposition_}\label{II:5-6}
$\tr_\Lambda(F) = |H|\tr_\Lambda^H(F)$.
\end{proposition_}

The formal properties of traces reduce us to treating only the
case where $P=\Lambda[H]$. The endomorphism $F$ is then right multiplication
by an element $\sum a_h\cdot h$ of $\Lambda[H]$, and
$\tr_\Lambda(F) = |H| a_e$, $\tr_\Lambda^H(F) = a_e$.




\subsection{}\label{II:5-7}

Let $A$ be a commutative ring and $\Lambda$ an $A$-algebra. Multiplication
by an element of $A$ passes to the quotient to define an $A$-module
structure on $A^\natural$. If $H$ is a finite group, $P$ an $A[H]$-module
projective and $M$ a $\Lambda[H]$-module, projective as a $\Lambda$-module,
one knows that the $\Lambda$-module $P\otimes_A H$, equipped with the diagonal action of
$H$, is a projective $\Lambda[H]$-module. Indeed:
\begin{enumerate}[\indent a)]
  \item It suffices to check it for $P=A[H]$.
  \item The $\Lambda$-module $P\otimes_A M$, equipped with the action of $H$ defined
    by its action on the first factor, is clearly a
    projective $\Lambda[H]$-module; denote it $(P\otimes_A M)'$.
  \item The map $h\otimes m\mapsto h\otimes h^{-1} m$ induces an
    isomorphism
    \begin{equation*}\tag{5.7.1}\label{II:eq:5-7-1}
      A[H]\otimes_A M \iso (A[H]\otimes_A M)' \text{.}
    \end{equation*}
\end{enumerate}




\begin{proposition_}\label{II:5-8}
With the previous notations, if $u$ is an endomorphism of $F$ and
$v$ an endomorphism of $M$, we have
\[
  \tr_\Lambda^H(u\otimes v) = \tr_A^H(u)\tr_\Lambda(v) \text{.}
\]
\end{proposition_}

We reduce to assuming that $P=A[H]$; the endomorphism $u$ is then right
multiplication $m_x$ by an element $x=\sum a_h h$ of $A[H]$.
Isomorphism \eqref{II:eq:5-7-1} transforms $u\otimes v$ into
$\sum a_h m_h\otimes h^{-1} v$, of trace $a_e\cdot\tr_\Lambda(v)$.




\subsection{}\label{II:5-9}

Let $G$ be a group extension of $\dZ$ by a finite group $H$
\[
  1 \to H \to G \to \dZ \to 1 \text{,}
\]
and $G^+$ the inverse image of $\dN$. Let $\Lambda$ be a ring, and $P$ a
$\Lambda[G^+]$-module, projective as a $\Lambda[H]$-module. If
$g\in G^+$, denote by $Z_g$ the centralizer of $g$ in $H$; multiplication
by $g$ is an endomorphism of $P$, seen as a $\Lambda[Z_g]$-module. $P$ being
$\Lambda[Z_g]$-projective, a trace $\tr_\Lambda^{Z_g}(g)$ is defined.
By \ref{II:5-6}, we have
\begin{equation*}\tag{5.9.1}\label{II:eq:5-9-1}
  \tr_\Lambda(g) = |Z_g| \cdot \tr_\Lambda^{Z_g}(g) \text{.}
\end{equation*}

Suppose that $\Lambda$ is an algebra over a commutative ring $A$. For
$P$ an $A[G^+]$-module, projective as an $A[H]$-module and $M$ a
$\Lambda[G^+]$-module, projective as a $\Lambda$-module, the
$\Lambda[G^+]$-module $P\otimes_A M$ (diagonal action of $G^+$) is projective
as a $\Lambda[H]$-module and by \ref{II:5-8}, for every
$g\in G^+$, we have
\begin{equation*}\tag{5.9.2}\label{II:eq:5-9-2}
  \tr_\Lambda^{Z_g}(g,P\otimes_A M) = \tr_A^{Z_g}(g,P) \cdot \tr_\Lambda(g,M) \text{.}
\end{equation*}




\subsection{}\label{II:5-10}

Let $P$ be a $\Lambda[G^+]$-module projective as a $\Lambda[H]$-module,
and $P_H$ the coinvariants of $H$ in $P$. It is a projective $\Lambda$-module.
The action of $g\in G^+$ passes to the quotient, and defines an endomorphism of
$P_H$ which depends only on the image of $g$ in $\dN$. We denote by $F$
the action of the elements $g\in G^+$ of image $1$.




\begin{proposition_}\label{II:5-11}
$\tr_\Lambda(F,P_H) = \sum_{g\mapsto 1}' \tr_\Lambda^{Z_g}(g,P)$, where
$\sum_{g\mapsto 1}'$ denotes a sum extended over the
$H$-conjugacy classes of elements of $G^+$ of image $1$ in $\dN$.
\end{proposition_}

After multiplication by $|H|$, this formula can be checked as follows:
\[
  |H|\cdot \tr_\Lambda(F,P_H)
    = \tr_\Lambda\left(\sum_{g\mapsto 1} g,P\right)
    = \sum_{g\mapsto 1}' \frac{|H|}{|Z_g|}\tr_\Lambda(g,P)
    = \sum_{g\mapsto 1}' |H|\cdot \tr_\Lambda^{Z_g}(g,P) \text{.}
\]

Choose an element $g\in G^+$ of image $1$ in $\dN$, and let
$\sigma$ be the automorphism of $\Lambda[H]$ induced by the automorphism
$\operatorname{ad} g$ of $H$. It amounts to the same to give a
$\Lambda[G^+]$-module, or a $\Lambda[H]$-module equipped with a
$\sigma$-linear endomorphism $\gamma$: one lets $g^n h$ act by $\gamma^n h$. In this
language, formula \ref{II:5-11} takes the form: for $P$ a
projective $\Lambda[H]$-module and for $\gamma$ a
$\sigma$-linear endomorphism of $P$, we have
\[
  \tr_\Lambda(\gamma,P_H) = \sum_h' \tr_\Lambda^{Z_g}(h\gamma,P)\text{,}
\]
where $\sum'$ is a sum over the $\operatorname{ad}g$-conjugacy classes
in $H$ (orbits of $H$ acting on itself by the
$x\mapsto h x \operatorname{ad} g(h)^{-1}$).

In this form, the properties of traces at once reduce us to the
case where $P=\Lambda[H]$. The endomorphism $Y$ then has the form
$x\mapsto \sigma(x) a$, with $a=\sum a_h h$ in $\Lambda[H]$. It suffices to
treat the universal case where the $a_h$ are indeterminates. In this
case, $\Lambda=\dZ\langle a_h\rangle_{h\in H}$ (noncommutative polynomials) and
$\Lambda^\natural$ is torsion free: one obtains \ref{II:5-11} by the argument
at the beginning and division by $|H|$. One can also compute directly.




\subsubsection{Remark}\label{II:5-11-1}

Everything said above also applies to
$\dZ/2$-graded modules, or to complexes of modules (cf. \ref{II:4-2}, \ref{II:4-3}).




\subsection{}\label{II:5-12}

Resume the notations of \ref{II:5-2}, and let $A=\dZ/\ell^n$, $n$ being
large enough for $\Lambda$ to be an $A$-algebra. Let $f:V_0\to U_0$ be a
finite \'etale connected Galois covering of $U_0$, with Galois group
$H$, and such that the sheaf $f^*\sG_0$ is constant. Denote by $M$ the constant
value $\h^0(V_0,f^*\sG_0)$, it is a $\Lambda[H]$-module, projective as
a $\Lambda$-module.

The sheaf $f_* A$ on $U_0$, equipped with the natural action of $H$, is a
locally free (of rank one) sheaf of $A[H]$-modules. The complex
$\eR\Gamma_c(U,f_* A)$ is hence an object of $\eD_\text{parf}\left(A[H]\right)$. It
is equipped with a Frobenius endomorphism $F^*$.

For the natural action of $H$ on $f_* f^*\sG_0$, the trace morphism
$f_* f^*\sG_0\to\sG_0$ factors through an isomorphism
\[
  \left(f_* f^*\sG_0\right)_H \to \sG_0 \text{.}
\]
Moreover, we have $f_* f^*\sG_0=f_* A\otimes_A M$ (with diagonal action of
$H$). The isomorphism
\[\xymatrix{
  \sG_0
    & \ar[l]_-\sim (f_* A\otimes_A M)_H = (f_* A\otimes_A M)\otimes_{\Lambda[H]} \Lambda
}\]
derives (\ref{II:4-12}) to an isomorphism
\[
  \eR\Gamma_c(U,\sG) \sim \left(\eR\Gamma_c(U,f_* A)\otimes_A M\right)\otimes_{\Lambda[H]} \Lambda \text{.}
\]

The endomorphism $F^*$ of the left hand side is deduced from the endomorphism
$F^*$ of $\eR\Gamma_c(U,f_* A)$. The complex of $\Lambda[H]$-modules
$\eR\Gamma_c(U,f_* A)$, equipped with $F^*$, can be seen as a complex of
$\Lambda[G^+]$-modules, for $G^+=H\times \dN$. Formulas \ref{II:5-10}
and \eqref{II:eq:5-9-2}, amplified by \ref{II:5-11-1}, hence provide
\[
  \tr\left(F^*,\eR\Gamma_c(U,\sG)\right) = \sum_h' \tr_A^{Z_g}\left(h F^*,\eR\Gamma_c(U,f_* A)\right)\cdot \tr_\Lambda(h,M)
\]
where $\sum'$ is a sum over the conjugacy classes in $H$, Moreover,
\[
  |Z_h|\cdot \tr_A^{Z_h}\left(h F^*,\eR\Gamma_c(U,f_ A)\right) = \tr_A\left((F h^{-1})^*,\eR\Gamma_c(V,A)\right) \text{.}
\]
This formula remains true for $A=\dZ/\ell^m$, with $m$ larger and larger.
Passing to the limit, one finds that
$\tr_A^{Z_g}\left(h F^*,\eR\Gamma_c(U,f^* A)\right)$ equals
\[
  \frac{1}{|Z_h|} \sum (-1)^i \tr\left((F h^{-1})^*,\h_c^i(V,\dQ_\ell)\right)
\]
(an element of $\dZ_\ell$), and
\begin{equation*}\tag{5.12.1}\label{II:eq:5-12-1}
\tr\left(F^*,\eR\Gamma_c(U,\sG)\right) = \sum_h' \frac{1}{|Z_h|} \tr\left((F h^{-1})^*,\h_c^i(V,\dQ_\ell)\right)\cdot \tr_\Lambda(h,M) \text{.}
\end{equation*}




\subsection{}\label{II:5-13}

It now remains to apply \ref{II:5-4} (note that if $\bar V$ is the
non-singular projective completion of $V$, the fixed points of $F h^{-1}$
are all of multiplicity one). One can do it without any computation:
\begin{enumerate}[\indent A.]
  \item If $U^F=\varnothing$, $\tr\left(F^*,\eR\Gamma_c(U,\sG)\right) = 0$. \end{enumerate}

Indeed, $F h^{-1}$ has no fixed point on $V$, \ref{II:5-4} shows that
$\sum (-1)^i \tr\left((F h^{-1})^*,\h_c^i(V,\dQ_\ell)\right)=0$ and one
applies \eqref{II:eq:5-12-1}.
\begin{enumerate}[\indent A.]
\setcounter{enumi}{1}
  \item In the general case, let $j:U'=U\setminus U^F\hookrightarrow U$.
    The filtration  $j_! j^* \sG_0\subset \sG_0$ of $\sG_0$ induces a
    filtration of $\eR\Gamma_c(U,\sG)$ with successive quotients
    $\eR\Gamma_c(U',\sG)$ and $\eR\Gamma_c(U^F,\sG)$. By \eqref{II:eq:4-4-1}
    and \ref{II:5-1}, we have
    \[
      \tr\left(F^*,\eR\Gamma_c(U,\sG)\right) = \tr\left(F^*,\eR\Gamma_c(U',\sG)\right) + \sum_{x\in U^F} \tr(F^*,\sG_x) \text{,}
    \]
    and, since ${U'}^F=0$, the first term on the right hand side is $0$.
\end{enumerate}











\section{End of the proof of \ref{II:4-10}}\label{II:6}

Prove \ref{II:4-10}, under the additional hypotheses that the ring
$\Lambda$ is killed by a power of a prime $\ell\ne p$, and
that $n=1$. The general case will follow: decompose $\Lambda$ into its
$\ell$-primary components ($\ell$ running over primes), and
replace $\dF_q$ by $\dF_{q^n}$, $X_0$ by $X_0\otimes_{\dF_q}\dF_{q^n}$.

For $X_0$ separated of finite type over $\dF_q$, and
$\sK_0\in\ob\eD_\text{ctf}(X_0,\Lambda)$, put
\[
  T'(X_0,\sK_0) = \sum_{x\in X^F} \tr(F^*,\sK_x)
\]
and
\[
  T''(X_0,\sK_0) = \tr(F^*,\eR\Gamma_c(X,\sK)) \text{.}
\]
One must prove the identity
\[
  (1;X_0,\sK_0) \qquad T'(X_0,\sK_0) = T''(X_0,\sK_0) \text{.}
\]
One begins by observing that $T'$ and $T''$ (generically denoted $T$)
obey the same formalism:

(A) If $j:U_0\hookrightarrow X_0$ is an open of $X_0$, with closed
complement $Y_0$, we have
\[
  T(X_0,\sK_0) = T(U_0,\sK_0|U_0) + T(Y_0,\sK_0|Y_0) \text{.}
\]
This is clear for $T'$. For $T''$, one observes that the filtration
$0\subset j_! j^* \sK_0\subset \sK_0$ of $\sK_0$ induces a filtration of
$\eR\Gamma_c(X,\sK)$, with successive quotients $\eR\Gamma_c(U,\sK)$ and
$\eR\Gamma_c(Y,\sK)$, and one invokes \eqref{II:eq:4-4-1} (more exactly
stated: $\sK_0$, equipped with the so-called filtration, defines an object of
$\mathsf{DF}(X,\Lambda)$; applying $\eR\Gamma_c$, one deduces from it an
object of $\mathsf{DF}_\text{parf}(\Lambda)$, to which $\eR\Gamma_c(X,\sK)$ is
underlying, and with graded as stated).

(B) If $\sK_0$ is a bounded complex with constructible flat components,
\[
  T(X_0,\sK_0) = \sum (-1)^i T(X_0,\sK^i) \text{.}
\]
This is clear for $T'$; for $T''$, one applies the above argument to the
stupid filtration of $\sK_0$, to obtain
\[
  T''(X_0,\sK_0) = \sum T''(X_0,\sK_0^i[-i]) = \sum (-1)^i T''(X_0,\sK_0^i) \text{.}
\]

(C) For $f:X_0\to Y_0$ a morphism, we have
\[
  T''(X_0,\sK_0) = T''(Y_0,\eR f_! \sK_0)
\]
and the same identity holds for $T'$ if the fibres of $f$ at the
rational points of $Y_0$ satisfy $\left(1;f^{-1}(y),\sK_0|f^{-1}(y)\right)$: one
uses that
\[
  T'(X_0,\sK_0) = \sum_{y\in Y^F} T'\left(f^{-1}(y), \sK_0|f^{-1}(y)\right)
\]
and
\[
  \eR\Gamma_c(Y, \eR f_! \sK) = \eR\Gamma_c(X,\sK) \text{.}
\]

Using this formalism, and \ref{II:4-7}, one reduces the proof of \ref{II:4-10}
to the two particular cases \ref{II:5-1} and \ref{II:5-2}.
